% ============================================================================
% Section 1: Introduction
% ============================================================================

\section{Introduction}
\label{sec:introduction}

\subsection{The hard problem of consciousness}

A scientific theory of consciousness must account for experience --- which is subjective ---
in objective terms. Being conscious, or having an experience, is understood to mean that
``there is something it is like to be'' in that state~\cite{nagel1974}: something it is like
to see a blue sky, hear the ocean roar, dream of a friend's face, or reflect on the experience
one is having.

In 1995, David Chalmers articulated the distinction between the ``easy problems'' of
consciousness (how the brain processes information, integrates sensory data, controls behavior)
and the ``hard problem'': why do these physical processes give rise to subjective experience
at all~\cite{chalmers1995}? The easy problems are engineering challenges --- technically
demanding but methodologically clear. The hard problem is of a different kind entirely:
even a complete functional description of the brain would leave unexplained \emph{why} there
is something it is like to undergo those functional processes.

Over 325 theories of consciousness exist~\cite{consciousness-atlas}, yet none has provided a
mathematical framework that simultaneously (a) addresses the explanatory gap between physics
and experience, (b) derives specific numerical thresholds for the onset of consciousness,
(c) connects consciousness to the laws of physics (quantum mechanics, general relativity),
and (d) makes falsifiable predictions with concrete experimental protocols.

\subsection{Why existing theories are insufficient}

The major existing theories each capture an important aspect of consciousness but fail to
provide a complete account:

\paragraph{Integrated Information Theory (IIT 4.0)~\cite{iit4}.}
IIT identifies five axioms of phenomenal existence (intrinsicality, information,
integration, exclusion, composition) and translates them into physical postulates.
Its central measure, $\Phi$ (integrated information), quantifies the irreducibility
of a system's cause--effect structure. However, IIT faces three critical
limitations: (1)~computing $\Phi$ requires searching over all partitions of the
system, whose count grows super-exponentially; practical full-$\Phi$ calculations
are limited to systems of roughly ten to fifteen elements, with approximation
schemes required beyond that; (2)~IIT provides no numerical threshold for
\emph{when} $\Phi$ becomes consciousness --- simple systems such as grids of
logic gates can have $\Phi > 0$ without clear consciousness; (3)~IIT does not
explain \emph{why} integrated information should feel like anything.

\paragraph{Global Neuronal Workspace Theory (GWT/GNWT)~\cite{dehaene2011}.}
GWT proposes that consciousness arises when information is broadcast to a ``global workspace''
of interconnected cortical neurons, enabling it to be accessed by multiple cognitive processes.
The theory successfully predicts the ``ignition'' pattern observed in EEG/fMRI studies.
However, GWT does not explain why broadcasting information to a workspace should produce
subjective experience (a loudspeaker broadcasts information without feeling anything), and
it makes no numerical predictions about the threshold of conscious access.

\paragraph{Free Energy Principle (FEP)~\cite{friston2019}.}
FEP provides a mathematical framework for understanding how self-organizing systems maintain
themselves by minimizing variational free energy. It elegantly unifies perception, action,
and learning under a single principle. However, FEP does not connect free energy minimization
to subjective experience --- a thermostat minimizes free energy without being conscious ---
and does not provide a threshold below which consciousness ceases.

\paragraph{Higher-Order Theories (HOT)~\cite{lau2011}.}
HOT proposes that a mental state becomes conscious when it is the object of a higher-order
representation --- roughly, when the system represents itself as being in that state.
This captures an important aspect of consciousness (self-awareness), but reduces consciousness
to meta-representation without explaining why meta-representation feels like anything, and
provides no quantitative predictions.

\paragraph{Predictive Processing (PP)~\cite{clark2013}.}
PP proposes that the brain operates as a hierarchical prediction machine, minimising
prediction error through generative models. Like FEP (with which it shares
mathematical foundations), PP does not explain why prediction should generate
subjective experience --- a prediction-error-minimising machine could operate
entirely unconsciously --- and provides no threshold for consciousness onset.

\paragraph{Conscious Realism / Conscious Agent Theory (CAT)~\cite{hoffman2014}.}
Hoffman and Prakash's Conscious Agent Theory is among the most formally
developed consciousness-first ontologies. Its Definition~1 specifies a
conscious agent as the 6-tuple
\[
  C \;=\; \bigl((X, \mathcal{X}),\; (G, \mathcal{G}),\; P,\; D,\; A,\; N\bigr),
\]
where $(X, \mathcal{X})$ is a measurable space of experiences,
$(G, \mathcal{G})$ a measurable space of actions, and $P$, $D$, $A$ are
Markovian kernels (perception, decision, action) linking these spaces to a
world $W$; $N$ is an integer counter tracking message passes. The paper proves
an \emph{Undirected Join Theorem} (Theorem~1) and a \emph{Directed Join
Theorem} (Theorem~2), establishing that conscious agents combine into larger
conscious agents. In Hoffman's broader programme, natural selection is argued
to favour fitness over veridical perception (the ``Fitness Beats Truth''
result~\cite{prakash2021fbt}), supporting the claim that spacetime is a
species-specific interface rather than fundamental reality.

CAT has three structural gaps relative to a complete theory:
(1)~\emph{no threshold for consciousness} --- every conscious agent is
conscious by definition, so CAT cannot distinguish a thermostat from a mammal
or say when a physical system \emph{becomes} a conscious agent;
(2)~\emph{no derivation of physics} --- spacetime is treated as illusory rather
than derived, leaving an unexplained gap to the empirical success of GR and QM;
(3)~\emph{no numerical predictions} --- the 2014 paper's content is structural
(join theorems, kernel composition), not quantitative, so it cannot be
falsified by any single measurement. UHM shares CAT's monistic starting point
but derives computable thresholds, emergent spacetime, and 22 numerical
predictions from the same primitive (see also \S\ref{sec:comparison}).

\paragraph{Projective Wave Theory (PWT)~\cite{worden2026}.}
Worden's Projective Wave Theory (2026) proposes that the brain's internal
model of local 3D space is held not in neurons but in a \emph{wave
excitation} carrying a projective transformation of Euclidean space, and that
this wave --- not its neural substrate --- is the source of spatial
conscious experience. Indirect evidence is adduced for such a wave in the
mammalian thalamus and in the central body of the insect brain. PWT is the
most recent serious competitor in the dynamical-structure family and agrees
with UHM that the substrate of experience is wave-like rather than
purely neural-computational. However, PWT targets a specific sub-problem
(the undistorted spatial form of consciousness) and a specific biological
substrate; it makes no numerical predictions, derives no physics,
formalises no threshold for when consciousness begins, and does not address
the hard problem directly. A systematic comparison is given in
\S\ref{sec:comparison}.

\medskip
In summary, each existing theory captures an important aspect of consciousness but none
provides a mathematical framework that simultaneously addresses the explanatory gap,
derives specific numerical thresholds, connects consciousness to the laws of physics,
and makes falsifiable numerical predictions with experimental protocols.

\subsection{Motivation: why a unified framework may be necessary}

The difficulty of resolving the hard problem within existing physics may
stem from a fundamental circularity inherent in the approach:

\begin{enumerate}[leftmargin=2em]
  \item To explain why neurons feel, one must explain what neurons \emph{are}.
  \item To explain what neurons are, one must explain matter.
  \item To explain matter, one needs quantum mechanics.
  \item To explain quantum mechanics, one must address measurement --- and in several
        interpretations (Copenhagen, von Neumann--Wigner), measurement is the point where
        an \emph{observer} enters the picture.
  \item In these interpretations, the observer is consciousness.
\end{enumerate}

The circle closes. Consciousness refers to physics, which refers to measurement, which
refers to consciousness. Every theory that explains consciousness within existing physics
inherits this circularity: ``information'' presupposes a subject who is informed;
``workspace'' presupposes someone working in it; ``free energy'' is defined relative to a
model, and a model requires a modeler.

This observation motivated a different approach: rather than explaining
consciousness \emph{within} physics, the present work asks whether
consciousness and physics might both be \emph{projections} of something more
fundamental. The resulting theory --- Unitary Holonomic Monism (UHM) ---
derives quantum mechanics, general relativity, and the structure of conscious
experience from a single mathematical primitive. The scope of the derivation
was not planned; it emerged from following the formalism to its logical
conclusions.

\subsection{Overview of UHM}

UHM proposes that reality is fundamentally described by a single object: the
\textbf{coherence matrix} $\Gam \in \Dens(\mathbb{C}^7)$ --- a $7 \times 7$ positive
semidefinite Hermitian matrix with unit trace. This object has two inseparable aspects:

\begin{itemize}[leftmargin=2em]
  \item \textbf{External aspect (physics):} The eigenvalues, dynamics, and symmetries of
        $\Gam$ encode physical structure --- from quantum states to spacetime geometry.
  \item \textbf{Internal aspect (experience):} The spectral decomposition of $\Gam$ on
        the E-dimension (interiority) encodes subjective content --- qualia, feelings,
        the sense of being.
\end{itemize}

These are not two different objects connected by a mysterious bridge. They
are one object seen from two angles --- like a coin's obverse and reverse.
This position, \textbf{two-aspect monism}, reframes the hard problem: there
is no explanatory gap because there was never a separation.

The theory is built on four axioms ($\Omega^{7}$, with a fifth derived from
the first four) grounded in the algebraic structure of octonions $\Oct$ and
the combinatorial structure of the Fano plane $\Fano$. Crucially, all four
axioms are either structurally forced or derived as theorems:
A1 (reality is an $\infty$-topos) is the minimal structure admitting internal
logic and differential cohesion~\cite{lurie2009}; A2 (Bures metric) is the
uniquely canonical Petz-enrichment, fixed by three independent characterizations --- Petz extremality, Uhlmann purification, SLD-Fisher Cram\'{e}r--Rao saturation (T-187~\status{T});
A3 ($N = 7$) follows from the Hurwitz--Adams classification of normed
division algebras together with Fano incidence; and A4 ($\omega_{0}$) is a
derived spectral property of the effective Hamiltonian $H_{\mathrm{eff}}$
(T-87~\status{T}). From these axioms, over 180 theorems follow, including:

\begin{itemize}[leftmargin=2em]
  \item A critical purity threshold $\Pcrit = 2/7 \approx 0.286$ below which the system
        irreversibly decays (Section~\ref{sec:theorems});
  \item A proof that philosophical zombies are dynamically impossible
        (No-Zombie Theorem, Section~\ref{subsec:no-zombie});
  \item Critical exponents of the consciousness phase transition:
        $\alpha = 1/2$, $\beta = 1/4$, $\gamma = 1$, $\nu = 1/2$, $\delta = 5$
        (Section~\ref{subsec:critical-exponents});
  \item Derivation of the Einstein field equations from $\Gam$
        (Appendix~\ref{app:einstein});
  \item 22 falsifiable numerical predictions with experimental protocols
        (Section~\ref{sec:predictions}).
\end{itemize}

The structure of the paper follows the logical chain of the theory:
Section~\ref{sec:axioms} presents the axioms and their motivation.
Section~\ref{sec:formalism} develops the mathematical formalism.
Section~\ref{sec:theorems} states and proves the key theorems.
Section~\ref{sec:cc} introduces Coherence Cybernetics --- the applied layer translating
formalism into measurable quantities (stress, hedonic valence, sensorimotor coupling).
Section~\ref{sec:predictions} collects all 22 predictions.
Section~\ref{sec:ai} derives implications for artificial consciousness and AGI.
Section~\ref{sec:ethics} examines the ethical consequences of the theorems.
Section~\ref{sec:comparison} compares UHM with existing theories.
Section~\ref{sec:experimental} outlines the experimental protocol.
Section~\ref{sec:discussion} synthesises the results and discusses limitations.
Appendices provide detailed proofs.

The complete formalization (192 theorems, T-1 through T-192, with full proofs)
constitutes the UHM theory archive. The present paper is self-contained for
its principal results: the appendices
prove $\Pcrit = 2/7$ (Appendix~\ref{app:pcrit}), the critical exponents
(Appendix~\ref{app:exponents}), and the Einstein-equation derivation
(Appendix~\ref{app:einstein}) without external dependencies.
